\documentclass[11pt]{article}

\usepackage[margin=1in]{geometry}
\usepackage{amsmath,amssymb,amsthm}
\usepackage{enumitem}
\usepackage{hyperref}
\usepackage{fancyhdr}
\usepackage{xcolor}
\usepackage{booktabs}

\hypersetup{colorlinks=true, linkcolor=blue!70!black, citecolor=blue!70!black, urlcolor=blue!70!black}

\pagestyle{fancy}
\fancyhf{}
\fancyhead[L]{\small Mathematical Review}
\fancyhead[R]{\small\thepage}
\renewcommand{\headrulewidth}{0.4pt}

\newcommand{\calA}{\mathcal{A}}
\newcommand{\calB}{\mathcal{B}}
\newcommand{\calF}{\mathcal{F}}
\newcommand{\calK}{\mathcal{K}}
\newcommand{\calX}{\mathcal{X}}
\newcommand{\RR}{\mathbb{R}}
\newcommand{\NN}{\mathbb{N}}
\newcommand{\inner}[2]{\left\langle #1,\, #2 \right\rangle}
\newcommand{\proj}[2]{P_{#1}\!\left[#2\right]}

\definecolor{issuecolor}{RGB}{180,40,40}
\definecolor{okcolor}{RGB}{30,120,50}
\definecolor{warncolor}{RGB}{200,130,0}

\newcommand{\verdict}[1]{\par\medskip\noindent\textbf{Verdict:} #1\par\medskip}
\newcommand{\statusok}{\textcolor{okcolor}{\textbf{[Correct]}}}
\newcommand{\statuswarn}{\textcolor{warncolor}{\textbf{[Minor issue]}}}
\newcommand{\statuserr}{\textcolor{issuecolor}{\textbf{[Needs correction]}}}

\theoremstyle{plain}
\newtheorem*{remarkstar}{Remark}

\title{\textbf{Review of the Mathematical Developments and Contributions}\\[6pt]
\large ``An augmented Lagrangian algorithm for constrained nonlinear least-squares''\\[4pt]
\normalsize P.~Borie, F.~Bastin}

\author{Detailed Technical Report}
\date{April 2026}

\begin{document}

\maketitle
\thispagestyle{fancy}

\begin{abstract}
This report presents a detailed analysis of the draft paper ``An augmented Lagrangian algorithm for constrained nonlinear least-squares.'' The first part (Sections~\ref{sec:overall}--\ref{sec:typos}) examines the validity of each major mathematical component: the augmented Lagrangian formulation, the projection machinery, the structured Hessian approximation, the Cauchy point computation, and the convergence theory. The second part (Section~\ref{sec:contributions}) assesses the three claimed contributions---the polyhedral inner solver, the structured SR1 update for the AL Hessian, and the convergence analysis---in terms of novelty, significance, and completeness. A summary of all findings is provided in Section~\ref{sec:summary}.
\end{abstract}

\tableofcontents

\newpage

%% ============================================================
\section{Overall Assessment}\label{sec:overall}
%% ============================================================

The paper is mathematically well-structured and the main arguments are sound. The augmented Lagrangian framework, the inner trust-region solver, and the convergence theory all follow established patterns from the optimization literature, primarily the work of Conn, Gould, and Toint. The paper makes a genuine contribution by combining the augmented Lagrangian framework with structured quasi-Newton updates for nonlinear least-squares in the polyhedral-constrained setting.

That said, the analysis below identifies several issues of varying severity, which we classify using the following labels:
\begin{itemize}[leftmargin=2cm]
  \item[\statusok] The development is mathematically correct.
  \item[\statuswarn] The development is essentially correct but contains a minor notational, presentational, or logical issue.
  \item[\statuserr] A substantive correction, additional justification, or missing proof is needed.
\end{itemize}

A summary table of all findings is provided in Section~\ref{sec:summary}.

%% ============================================================
\section{Problem Formulation and Augmented Lagrangian Setup}
%% ============================================================

\subsection{Problem Statement (Equations~(1.1)--(1.2))}
\statusok

The constrained nonlinear least-squares problem is stated as
\[
\min_{x\in\RR^n} \;\tfrac{1}{2}\|r(x)\|^2 \quad \text{s.t.}\quad c(x)=0,\quad x\in\Omega,
\]
with $\Omega = \{x \in \RR^n \mid Ax = b,\; l \le x \le u\}$. The assumptions that $r$ and $c$ are $C^2$, that $A$ has full row rank, and that $\Omega$ is non-empty are standard and sufficient for the subsequent developments.

\subsection{Lagrangian and First-Order Conditions (Equations~(1.3)--(1.5))}
\statuswarn

The Lagrangian $\mathcal{L}(x,\lambda) = f(x) + \inner{\lambda}{c(x)}$ is standard. The first-order condition for a critical point $(x^*,\lambda^*)$ is correctly stated as
\[
x^*\in\Omega,\quad c(x^*)=0,\quad \proj{\Omega}{x^*-\nabla_x\mathcal{L}(x^*,\lambda^*)}=x^*.
\]
This is the standard projected-gradient stationarity condition for minimization over a convex set and is correct as a \emph{necessary} condition for local optimality.

\subsection{Augmented Lagrangian Function (Equations~(2.1)--(2.5))}
\statusok

The AL function
\[
\Phi(x,\lambda,\mu) = \tfrac{1}{2}\|r(x)\|^2 + \inner{\lambda}{c(x)} + \tfrac{\mu}{2}\|c(x)\|^2
\]
is standard. The gradient
\[
\nabla_x\Phi = J(x)^Tr(x) + C(x)^T\bar{\lambda}(x,\lambda,\mu),
\]
with $\bar{\lambda} = \lambda + \mu\,c(x)$, is obtained by straightforward differentiation and is correct. The Hessian
\[
\nabla^2_{xx}\Phi = J^TJ + \mu\,C^TC + S(x),
\]
where $S(x) = \sum_{i=1}^{n_r} r_i\nabla^2 r_i + \sum_{i=1}^{n_c} \bar{\lambda}_i\nabla^2 c_i$, is verified by differentiating the gradient expression. Each term arises correctly from the product rule applied to $J^Tr$ and $C^T\bar{\lambda}$.

\subsection{Stationarity of the AL Subproblem (Equation~(2.6))}
\statuswarn

The paper writes
\[
x^*\in \arg\min_{x\in\Omega}\Phi(x,\lambda,\mu) \;\iff\; \proj{\Omega}{x^*-\nabla_x\Phi(x^*,\lambda,\mu)}=x^*.
\]
The right-to-left implication (``$\Leftarrow$'') is problematic: the projected-gradient condition only ensures first-order stationarity, not global minimality. That is, the condition $\proj{\Omega}{x^* - \nabla_x\Phi} = x^*$ is \emph{necessary} for $x^*$ to be a local minimizer of $\Phi$ over $\Omega$, but it is not sufficient (it could be a saddle point or a local maximum on the boundary). The ``$\arg\min$'' and ``$\iff$'' overstate what is actually established.

\verdict{Replace ``$x^*\in\arg\min$'' with ``$x^*$ is a first-order critical point of $\min_{x\in\Omega}\Phi$'' and change the biconditional to a necessary condition, or restrict the equivalence to the convex case.}


%% ============================================================
\section{Projection Machinery}
%% ============================================================

\subsection{Tangent Space and Reduced Gradient (Equations~(2.7)--(2.9))}
\statusok

The active set $I(x)$, the tangent space
\[
T(x) = \{d \in \RR^n \mid Ad = 0,\; d_i = 0\;\forall i\in I(x)\},
\]
and the reduced gradient $\proj{T(x)}{\nabla_x\Phi}$ are standard constructions for linearly constrained optimization. The discussion of its role as a criticality measure is appropriate and consistent with the literature (MINOS, LSNNO).

\subsection{Normal Equations Approach (Equations~(2.11)--(2.14))}
\statusok

The augmented constraint matrix $\tilde{A} = \begin{pmatrix} A \\ Z \end{pmatrix}$ where $Z$ selects the rows of the identity corresponding to active bounds is correctly formed. The projection formula
\[
\tilde{P} = I - \tilde{A}^T(\tilde{A}\tilde{A}^T)^{-1}\tilde{A}
\]
is the standard orthogonal projector onto the null space of $\tilde{A}$, valid when $\tilde{A}$ has full row rank. The computation via the normal equations---solve $(\tilde{A}\tilde{A}^T)y = \tilde{A}v$, then form $v - \tilde{A}^Ty$---is correct.

\medskip
\noindent\textcolor{warncolor}{\textbf{Remark on full rank of $\tilde{A}$.}} The claim that $\tilde{A}$ is full rank under Assumption~2 ($A$ full row rank) and $p < n-m$ deserves a brief justification. The rows of $Z$ are coordinate vectors $e_{i_1},\ldots,e_{i_p}$. For $\tilde{A}$ to be full rank, these must be linearly independent of the rows of $A$. This holds generically, but could fail in the degenerate case where a row of $A$ is itself a coordinate vector $e_{i_k}$ for some active bound index $i_k$. The paper should either add a non-degeneracy assumption or note that this situation cannot arise in practice (for instance, if the linear equalities do not involve single-variable constraints).

\subsection{Redundancy in the Projection Formula}
\statuswarn

The paper first introduces $\tilde{N}$ as a matrix whose columns form an \emph{orthonormal} basis of the null space of $\tilde{A}$, then writes the projector as $\tilde{P} = \tilde{N}(\tilde{N}^T\tilde{N})^{-1}\tilde{N}^T$. Since $\tilde{N}$ is orthonormal by construction, $\tilde{N}^T\tilde{N} = I$ and the formula simplifies to $\tilde{P} = \tilde{N}\tilde{N}^T$. The general expression is correct but unnecessarily complex given the stated orthonormality. This is a harmless redundancy, not an error.

\subsection{Block Cholesky Factorization (Appendix)}
\statuswarn

The block structure of $\tilde{A}\tilde{A}^T$ is correctly identified:
\[
\tilde{A}\tilde{A}^T = \begin{bmatrix} AA^T & A_{|\calA} \\ A_{|\calA}^T & I_p \end{bmatrix},
\]
where $A_{|\calA}$ denotes the restriction of $A$ to the columns indexed by the active set. The block Cholesky factorization $\tilde{L}\tilde{L}^T$ with
\[
\tilde{L} = \begin{bmatrix} \tilde{L}_{11} & 0 \\ G_{21} & \tilde{L}_{22} \end{bmatrix}
\]
is derived correctly: $\tilde{L}_{11} = L$ (the Cholesky factor of $AA^T$), $G_{21}$ is obtained by solving $LG_{21}^T = A_{|\calA}$, and $\tilde{L}_{22}$ is the Cholesky factor of $I_p - G_{21}G_{21}^T$.

\medskip\noindent
\textcolor{warncolor}{\textbf{Missing justification:}} The paper does not verify that $I_p - G_{21}G_{21}^T$ is positive definite. This property is necessary for the Cholesky factorization to exist. It follows from the positive definiteness of $\tilde{A}\tilde{A}^T$ (which in turn follows from the full-rank assumption on $\tilde{A}$), since the Schur complement of the $(1,1)$ block in a positive definite matrix is itself positive definite:
\[
I_p - A_{|\calA}^T(AA^T)^{-1}A_{|\calA} = I_p - G_{21}G_{21}^T \succ 0.
\]
This argument should be stated explicitly.

%% ============================================================
\section{Inner Minimization and Step Computation}
%% ============================================================

\subsection{Trust Region Subproblem (Equations~(2.15)--(2.19))}
\statusok

The quadratic model $q_k(s) = \frac{1}{2}\inner{s}{H_ks} + \inner{g_k}{s}$ and the trust region subproblem
\[
\min_s\; q_k(s) \quad\text{s.t.}\quad x_k + s \in \Omega,\;\|s\|_\infty \le \Delta_k
\]
are standard. The trust region ratio $\rho_k$ and the step acceptance/radius update criteria are correctly formulated and consistent with the referenced literature.

\subsection{Minor Iterations and Projected CG (Section~2.2)}
\statusok

The decomposition of the step into a Cauchy phase followed by minor iterations using the projected conjugate gradient (PCG) method is a well-established approach. The conditions on the minor iterates---feasibility, containment in the trust region, and nested active sets (Equation~(2.22))---are correctly stated. The sufficient decrease condition (Equation~(2.23)) between successive minor iterates and the stopping criterion based on the reduced gradient (Equation~(2.26)) are appropriate.

The PCG algorithm (Algorithm~3) correctly handles the three termination cases: bound hit (scale direction to the boundary), negative curvature (scale direction to the feasible boundary), and convergence (reduced gradient tolerance met). The iteration limit of $2(n - m - |I(x_{k,j})|)$ is reasonable.

\subsection{Magical Step for Slack Variables (Equations~(2.36)--(2.42))}
\statusok

The closed-form minimizer for the slack variables,
\[
(\hat{\nu}_k)_i = \max\!\left(0,\;\frac{\lambda_i}{\mu} + g_i(\bar{x}_k)\right),
\]
is correct. Since $\Phi$ is quadratic and convex in~$\nu$, this is obtained by setting $\partial\Phi/\partial\nu_i = 0$ and projecting onto $\nu_i \ge 0$. The modified predicted reduction (Equation~(2.42)), which combines the model-based prediction for the $x$-step with the actual reduction from the slack update, is a well-motivated and established approach.

%% ============================================================
\section{Structured Hessian Approximation}
%% ============================================================

\subsection{Gauss--Newton Approximation (Equation~(2.29))}
\statusok

The GN approximation $H_k = J_k^TJ_k + \mu C_k^TC_k$ is the natural extension of the unconstrained GN Hessian to the AL setting. The discussion of its limitations---neglecting both residual curvature and multiplier contributions---is appropriate.

\subsection{Structured Secant Equation (Equations~(2.30)--(2.34))}
\statuswarn

The derivation of the structured secant equation is the key methodological contribution. Starting from the decomposition $B_{k+1} = \sum_i r_i(x_{k+1})B_{k+1}^{r_i} + \sum_i \bar{\lambda}_i(x_{k+1})B_{k+1}^{c_i}$, imposing individual secant equations
\[
B_{k+1}^{r_i}s_k = \nabla r_i(x_{k+1}) - \nabla r_i(x_k), \qquad B_{k+1}^{c_i}s_k = \nabla c_i(x_{k+1}) - \nabla c_i(x_k),
\]
and summing yields the structured secant equation
\[
B_{k+1}s_k = (J_{k+1} - J_k)^Tr_{k+1} + (C_{k+1} - C_k)^T\bar{\lambda}_{k+1} =: y_k^A.
\]
This algebraic derivation is correct. We note two points.

\paragraph{Typo.} In the sentence following Equation~(2.33), the text writes ``$\nabla c_i(c_{k+1}) - \nabla c_i(x_k)$'', where ``$c_{k+1}$'' should be ``$x_{k+1}$''.

\paragraph{On the relationship between $y_k^A$ and $y_k$.} The structured right-hand side $y_k^A$ differs from the standard quasi-Newton difference $y_k = g_{k+1} - g_k$. Specifically:
\begin{align*}
y_k &= g_{k+1} - g_k = J_{k+1}^Tr_{k+1} + C_{k+1}^T\bar{\lambda}_{k+1} - J_k^Tr_k - C_k^T\bar{\lambda}_k, \\
y_k^A &= (J_{k+1} - J_k)^Tr_{k+1} + (C_{k+1} - C_k)^T\bar{\lambda}_{k+1}.
\end{align*}
Thus $y_k - y_k^A = J_k^T(r_{k+1} - r_k) + C_k^T(\bar{\lambda}_{k+1} - \bar{\lambda}_k) = B_k^{GN}s_k + O(\|s_k\|^2)$ to first order, which is precisely the GN contribution. This is by design: $y_k^A$ captures only the second-order part, and the full Hessian approximation is $H_{k+1} = B_{k+1}^{GN} + B_{k+1}$, where $B_{k+1}$ satisfies $B_{k+1}s_k = y_k^A$. The approach is consistent and mathematically sound, though it would benefit from a brief remark making this relationship explicit.

\subsection{SR1 Update Formula (Equations~(2.35)--(2.37))}
\statusok

The SR1 update
\[
B_{k+1} = B_k + \frac{(y_k^A - B_ks_k)(y_k^A - B_ks_k)^T}{\inner{y_k^A - B_ks_k}{s_k}}
\]
is the standard rank-one symmetric update applied with the structured right-hand side. The safeguard condition $|\inner{y_k^A - B_ks_k}{s_k}| \ge \kappa_{sds}\|s_k\|\|y_k^A - B_ks_k\|$ is a well-known numerical safeguard for SR1 methods. The discussion of why SR1 is preferred over BFGS/DFP in this context (no curvature condition needed, better Hessian approximation properties, compatibility with trust-region CG) is well-argued.


%% ============================================================
\section{Convergence Analysis: Inner Minimization}
%% ============================================================

\subsection{Setup and Assumptions}
\statusok

Assumptions~4 (compactness of the level set), 5 (divergence of $\sum 1/b_k$), and 6 ($\lim b_k(\varphi_{k+1} - \varphi_k) = 0$) are the standard conditions for trust-region global convergence with possibly unbounded Hessian approximations. The criticality measure $h_k = \|s_k(1)\|$ and the projected gradient path formalism are correctly set up.

\subsection{Projected Gradient Path Analysis}

\subsubsection{Lemma~3.1 (Non-increasing reduced gradient norms)}
\statusok

The two claims are: (i) $T(t) \supseteq T(t')$ for $t < t'$ and (ii) $\|z_k(t)\| \ge \|z_k(t')\|$. The first follows from the nested active sets $I(t) \subseteq I(t')$, which constrains the tangent space further. The second follows from the variational characterization of orthogonal projection: projecting the same vector onto a smaller subspace cannot increase the norm. Both are correct.

\subsubsection{Integral Representation (Equation~(3.14))}
\statusok

The expression
\[
s_k(t) = -\int_0^t z_k(\tau)\,d\tau,
\]
understood as a sum of integrals over the intervals $[t_j,t_{j+1}]$ where $z_k$ is constant, is a valid reformulation of the piecewise-linear recursion (3.10). Since $s_k'(t) = -z_k(t)$ on each open interval, and $s_k(0) = 0$, the fundamental theorem of calculus gives the result.

\subsubsection{Lemma~3.2 (Lower bound on reduced gradient norm)}
\statuserr

The lemma claims that $\|z_k(t_k^{(1)})\| \ge h_k/2$ where $t_k^{(1)} = h_k/(2\kappa_{ubg})$. The proof strategy is correct but the presentation of the chain of inequalities is flawed.

\paragraph{The non-expansiveness bound.} The inequality $\|s_k(t)\| \le t\|g_k\|$ is obtained via
\[
\|s_k(t)\| = \|\proj{\Omega}{x_k - tg_k} - \proj{\Omega}{x_k}\| \le \|tg_k\| = t\|g_k\|,
\]
using the non-expansiveness (1-Lipschitz property) of the projection onto the convex set $\Omega$. This is correct.

\paragraph{The problematic chain.} The proof then writes (around lines~822--829):
\[
\|z_k(t_k^{(1)})\| \ge (t_k^{(2)} - t_k^{(1)})\|z_k(t_k^{(1)})\| \ge \|s_k(t_k^{(2)}) - s_k(t_k^{(1)})\| \ge \|s_k(t_k^{(2)})\| - \|s_k(t_k^{(1)})\| \ge \frac{h_k}{2}.
\]
The first inequality requires $t_k^{(2)} - t_k^{(1)} \le 1$, which holds since both scalars are in $[0,1]$, but this ``$a \ge ba$'' step is logically correct only because $b \le 1$ and $a \ge 0$. More fundamentally, the chain is written in a confusing direction: the natural flow of the argument goes from right to left (from the known bound $h_k/2$ upward), but the inequalities are displayed left to right.

\verdict{The final result $\|z_k(t_k^{(1)})\| \ge h_k/2$ is correct. However, the proof should be rewritten with the chain of inequalities flowing in a single, natural direction. The clean version is:
\begin{align*}
\frac{h_k}{2} &= \|s_k(t_k^{(2)})\| - \frac{h_k}{2} \le \|s_k(t_k^{(2)})\| - \|s_k(t_k^{(1)})\| \\
&\le \|s_k(t_k^{(2)}) - s_k(t_k^{(1)})\|  \le (t_k^{(2)} - t_k^{(1)})\|z_k(t_k^{(1)})\| \le \|z_k(t_k^{(1)})\|,
\end{align*}
using the reverse triangle inequality, the integral bound~(3.19), and $t_k^{(2)} - t_k^{(1)} \le 1$.}

\subsection{Lemma~3.3 (Model Decrease Along the Projected Gradient Path)}
\statusok

The bound on the derivative of the piecewise quadratic $\psi(t) = q_k(s_k(t))$ is derived by splitting into first-order and second-order contributions:
\[
\psi'(t) = -\inner{g_k}{z_i} - \inner{s_k(t)}{H_kz_i}.
\]

\paragraph{First-order term.} Since $z_i = \proj{T_i}{g_k}$, one has $\inner{g_k}{z_i} = \|z_i\|^2 \ge \|z_k(t_k^{(3)})\|^2 = \alpha_k^2$ by the monotonicity of the reduced gradient norm (Lemma~3.1). This is correct.

\paragraph{Second-order term.} By Cauchy--Schwarz and the triangle inequality applied to the integral representation:
\[
|\inner{s_k(t)}{H_kz_i}| \le \|s_k(t)\|\cdot\|H_k\|\cdot\|z_i\| \le t_k^{(3)}\kappa_{ubg}^2\|H_k\|,
\]
using $\|s_k(t)\| \le t\|g_k\| \le t_k^{(3)}\kappa_{ubg}$ and $\|z_i\| \le \|g_k\| \le \kappa_{ubg}$. This is correct.

\paragraph{Definition of $t_k^{(4)}$.} The choice
\[
t_k^{(4)} = \min\!\left(t_k^{(3)},\;\frac{\alpha_k^2}{2\kappa_{ubg}^2(\|H_k\|+1)}\right)
\]
ensures $\psi'(t) \le -\alpha_k^2/2$ for $t \le t_k^{(4)}$, and hence $\psi(t) \le -\alpha_k^2 t/2$ by integration. The ``$+1$'' in the denominator prevents division by zero when $H_k = 0$. This is a standard safeguard. The derivation is correct.

\subsection{Lemma~3.4 (Guaranteed Decrease by the Step)}\label{subsec:lemma34}
\statuserr

The lemma establishes the key decrease bound
\[
q_k(s_k) \le -\kappa_{mdc}\,h_k^2\min\!\left(\frac{h_k^2}{b_k},\;\Delta_k\right), \qquad \kappa_{mdc} = \frac{\kappa_{fcd}}{64\kappa_{ubg}^2}.
\]
The proof splits into two cases depending on whether the Cauchy step at $t_k^{(5)}$ lies inside or at the boundary of the trust region.

\subsubsection{Case 1: Cauchy step inside the trust region}
\statuserr

The proof writes: ``assume that $\|s_k(t_k^{(5)})\|_\infty \le \Delta_k$. Then$\ldots$'' and directly invokes the fraction-of-Cauchy-decrease condition to conclude. However, the argument requires a justification step: since $\psi'(t) \le -h_k^2/8 < 0$ for all $t \in [0, t_k^{(5)}]$ (by Lemma~3.3 applied with $\alpha_k = h_k/2$), the function $\psi$ is strictly decreasing on this interval. Because the Cauchy point $t_k^C$ is the first local minimizer of $\psi$ subject to the trust region, and the trust region is not active at $t_k^{(5)}$, we must have $t_k^C \ge t_k^{(5)}$. Therefore $\psi(t_k^C) \le \psi(t_k^{(5)}) \le -h_k^2 t_k^{(5)}/8$, and the fraction-of-Cauchy-decrease condition gives the result. This gap corresponds to the TODO note at line~933 of the draft.

\verdict{Fill in the justification: the monotonicity of $\psi$ on $[0,t_k^{(5)}]$ implies $t_k^C \ge t_k^{(5)}$, hence $q_k(s_k^C) \le \psi(t_k^{(5)})$.}

\subsubsection{Case 2: Cauchy step at the trust region boundary}
\statusok

When $\|s_k(t_k^{(5)})\|_\infty > \Delta_k$, the proof shows $t_k^C \ge \Delta_k/\kappa_{ubg}$ using non-expansiveness, and bounds $q_k(s_k^C) \le -\kappa_{fcd}h_k^2\Delta_k/(8\kappa_{ubg})$. The final simplification to $-\kappa_{fcd}h_k^2\Delta_k/(64\kappa_{ubg}^2)$ using $\kappa_{ubg} \ge 1$ is correct.

\subsubsection{Verification of the constant $\kappa_{mdc}$}
\statusok

Combining both cases:
\begin{itemize}
  \item Case~1: $q_k(s_k) \le -\kappa_{fcd}\frac{h_k^2}{8}\cdot\frac{h_k^2}{8\kappa_{ubg}^2(\|H_k\|+1)} \le -\frac{\kappa_{fcd}}{64\kappa_{ubg}^2}\cdot\frac{h_k^4}{b_k}$.
  \item Case~2: $q_k(s_k) \le -\frac{\kappa_{fcd}}{64\kappa_{ubg}^2}\cdot h_k^2\Delta_k$.
\end{itemize}
Here we used $b_k = 1 + \max_{0\le i\le k}\|H_i\| \ge \|H_k\| + 1$. Taking the minimum gives $\kappa_{mdc} = \kappa_{fcd}/(64\kappa_{ubg}^2)$. Correct.

\subsubsection{Successful step inequality}
\statusok

For a successful step ($\rho_k > \eta_1$), the actual reduction satisfies
\[
\varphi(x_k) - \varphi(x_k + s_k) \ge \eta_1\,|q_k(s_k)| \ge \kappa_{sdc}\,h_k^2\min\!\left(\frac{h_k^2}{b_k},\Delta_k\right),
\]
with $\kappa_{sdc} = \eta_1\kappa_{mdc}$. This follows directly from the definition of $\rho_k$ and the sign of $q_k(s_k) < 0$. Correct.

\subsection{Main Convergence Theorem (Theorem~3.1)}
\statusok

The conclusion $\lim_{k\to\infty} h_k = 0$ is the standard global convergence result for trust-region methods with the given assumptions. The paper correctly notes that once the guaranteed decrease inequality (Lemma~3.4) is established, the remaining argument is standard and does not involve the polyhedral constraint structure. Referencing the proof from Conn et al.\ (1988a) is appropriate.

%% ============================================================
\section{Convergence Analysis: Outer Algorithm}
%% ============================================================

\subsection{Problem Reformulation}
\statusok

The observation that linear constraints $Ax = b$, $l \le x \le u$ can be written as $\bar{A}x \ge \bar{b}$ with the block structure is trivially correct and places the problem within the framework of Conn et al.\ (1996b).

\subsection{Assumptions~7--8}
\statusok

Assumption~7 (bounded iterates) is standard for global convergence of AL methods. Assumption~8 (rank of $C(x^*)N_*$ is at least $n_c$) is the appropriate constraint qualification: it ensures that the nonlinear constraint gradients are linearly independent within the null space of the active linear constraints at the limit point. This is a natural adaptation of LICQ to the polyhedral-constrained setting.

\subsection{Theorem~3.2 (Global Convergence)}
\statuserr

The theorem states that under Assumptions~1--3 and~7--8, any convergent subsequence of outer iterates converges to a feasible first-order critical point with corresponding multipliers. The statement is consistent with the referenced theory.

\paragraph{Missing proof.} The proof is entirely absent. The draft contains only the comment ``A sketch of the proof could be useful.'' While the result is stated as following from Conn et al.\ (1996b), the adaptation to the present setting (mixed equalities and bounds, AL function with least-squares structure, structured Hessian approximation) is non-trivial and requires verification of several intermediate steps:

\begin{enumerate}[label=(\roman*)]
  \item That the inner solver satisfies the requirements of the outer AL framework (guaranteed approximate stationarity at each outer iteration).
  \item That the multiplier update $\lambda_{k+1} = \bar{\lambda}(x_k,\lambda_k,\mu_k)$ and the penalty parameter update ensure the claimed asymptotic behavior.
  \item That the constraint qualification (Assumption~8) is sufficient to guarantee feasibility of the limit point, via a standard argument showing that unbounded penalty parameters imply constraint infeasibility at the limit, contradicting the CQ.
  \item That the least-squares multiplier estimate~(3.35) converges to the correct multiplier.
\end{enumerate}

\verdict{At minimum, a sketch of the proof should be provided, identifying which results from Conn et al.\ (1996b) are invoked and verifying that their hypotheses are met. This is the most significant gap in the current manuscript.}

%% ============================================================
\section{Cauchy Point Computation (Appendix)}
%% ============================================================

\subsection{Projected Gradient Path with Linear Equalities}
\statusok

The adaptation of the Cauchy point computation to the polyhedral case is correctly presented. The path $s(t) = \proj{\Omega}{x - tg} - x$ satisfies $As(t) = 0$ and the component-wise bounds $s^{(l)} \le s(t) \le s^{(u)}$, where $s^{(l)}_i = \max(-\Delta, l_i - x_i)$ and $s^{(u)}_i = \min(\Delta, u_i - x_i)$. The recursive expression $s(t) = -(t - t_i)z_i + s(t_i)$ for $t \in [t_i, t_{i+1})$ is consistent with the reduced gradient $z_i$ defined in the main text.

\subsection{Case Analysis for the Quadratic Model}
\statusok

The piecewise quadratic $\psi(t) = q(s(t))$ on each interval $[t_i, t_{i+1})$ has the form $\psi(t) = \frac{\psi_i''}{2}(\Delta t)^2 + \psi_i'\Delta t$ with $\Delta t = t - t_i$, $\psi_i'' = \inner{z_i}{Hz_i}$, and $\psi_i' = -\inner{s(t_i)}{Hz_i} - \inner{g}{z_i}$. The case analysis is standard and correct:

\begin{itemize}
  \item $\psi_i' > 0$ or ($\psi_i' = 0$ and $\psi_i'' > 0$): $t_i$ is the minimizer.
  \item $\psi_i' < 0$ and $\psi_i'' > 0$: the minimizer is at $t_i - \psi_i'/\psi_i''$ if it lies in $[t_i, t_{i+1})$.
  \item Otherwise: proceed to the next interval.
\end{itemize}

\subsection{Breakpoint Computation}
\statusok

The formulas for $\delta_i^-$ and $\delta_i^+$, which identify the next bound hit along the direction $-z_i$, are correct. The termination condition (all $n - m$ bounds active) correctly accounts for the $m$ degrees of freedom consumed by the equality constraints.

\subsection{Algorithm~4 (Cauchy Step Computation)}
\statusok

The algorithm correctly implements the case analysis and breakpoint computation. The loop invariant---maintaining the current step $s^{(i)}$, the breakpoint $t_i$, and the reduced gradient $z_i$---is consistent with the mathematical description.

%% ============================================================
\section{Tolerance Update Scheme (Algorithm~1)}
%% ============================================================
\statuswarn

The tolerance updates in Algorithm~1 follow two distinct paths depending on whether sufficient feasibility progress is achieved:

\begin{itemize}
  \item \textbf{Feasibility improved} ($\|c(x_k)\| \le \eta_k$): update multipliers, keep $\mu$, and tighten tolerances via $\omega_{k+1} = \omega_k\mu_{k+1}^{-\beta_\omega}$, $\eta_{k+1} = \eta_k\mu_{k+1}^{-\beta_\eta}$.
  \item \textbf{Feasibility not improved}: keep $(x,\lambda)$, increase $\mu_{k+1} = \tau\mu_k$, and reset tolerances via $\omega_{k+1} = \omega\mu_{k+1}^{-\kappa_\omega}$, $\eta_{k+1} = \eta\mu_{k+1}^{-\kappa_\eta}$.
\end{itemize}

This scheme is consistent with the standard AL update rules from Conn et al.\ (2000, Chapter~14). Since $\mu_k \to \infty$ whenever feasibility is not eventually achieved, the tolerances $\omega_k,\eta_k \to 0$, which is required for convergence.

\medskip\noindent
\textcolor{warncolor}{\textbf{Minor observation.}} In the ``feasibility improved'' branch, the update $\omega_{k+1} = \omega_k\mu_{k+1}^{-\beta_\omega}$ with $\mu_{k+1} = \mu_k$ (unchanged penalty) gives $\omega_{k+1} = \omega_k\mu_k^{-\beta_\omega}$. For the tolerances to converge to zero even when $\mu$ stays fixed (because feasibility keeps improving), we need $\beta_\omega > 0$, which is the case with the default $\beta_\omega = 1$. However, the convergence to zero is only guaranteed if the number of ``feasibility improved'' iterations is infinite, or if $\mu_k\to\infty$. This is a standard subtlety of AL methods and is handled in the referenced theory.


%% ============================================================
\section{Minor Notational and Typographical Issues}\label{sec:typos}
%% ============================================================

Beyond the substantive issues discussed above, we note the following minor points:

\begin{enumerate}[label=\arabic*.]
  \item \textbf{Equation~(2.33), line~521:} ``$\nabla c_i(c_{k+1}) - \nabla c_i(x_k)$'' should be ``$\nabla c_i(x_{k+1}) - \nabla c_i(x_k)$''.
  
  \item \textbf{Equation~(2.6):} The biconditional ``$\iff$'' should be weakened to ``$\implies$'' (right-to-left) or the left-hand side should be changed to ``first-order critical point'' instead of ``$\arg\min$'' (see Section~2.4 above).
  
  \item \textbf{Line~625:} The projected gradient norm is written as $\|x_K - \proj{\Omega}{x_K - \Phi_K}\|$, where ``$\Phi_K$'' should be ``$\nabla_x\Phi_K$''.

  \item \textbf{Lemma~3.2, statement:} The constant $\kappa_{ubg}$ appears in the lemma statement (line~771) but is not used in the conclusion ``$\|z_k(t_k^{(1)})\| \ge h_k/2$''. It is only used to define $t_k^{(1)} = h_k/(2\kappa_{ubg})$. The statement could be streamlined.
  
  \item \textbf{Lemma~3.4, proof, line~928:} The notation ``$t_1^{(k)}$'' appears to be a reversal of the subscript/superscript convention; it should be ``$t_k^{(1)}$''.

  \item \textbf{Section~2.2, Equation~(2.26):} The tangent space is written $T_{k,j}$ but defined ``at $x_{x,j}$'' (double subscript typo), which should be ``at $x_{k,j}$''.
\end{enumerate}


%% ============================================================
\section{Assessment of the Contributions}\label{sec:contributions}
%% ============================================================

The paper claims three contributions (Introduction, paragraph~5): (1)~an AL algorithm with an inner solver designed for polyhedral feasible sets, (2)~a structured SR1 update exploiting the least-squares structure of the AL Hessian, and (3)~a global convergence analysis. We assess each in turn, then discuss what is missing.

\subsection{Contribution 1: AL Framework with Polyhedral Inner Solver}

\paragraph{Claim.} The inner minimization phase is specifically designed for polyhedral feasible sets $\Omega = \{x \mid Ax = b,\; l \le x \le u\}$, combining bound constraints with linear equalities. The inner solver uses a gradient projection technique (Cauchy point along the projected gradient path, followed by minor iterations via projected CG), with efficient projections computed through a normal equations approach with an incrementally updated block Cholesky factorization.

\paragraph{Assessment.} This is the most concrete and practically significant contribution. The LANCELOT solver, which is the closest algorithmic relative, handles only bound constraints in its AL subproblems. Extending to the mixed bounds-plus-linear-equalities case requires non-trivial machinery:
\begin{itemize}
  \item The Cauchy point computation along the projected gradient path must account for the fact that ``hitting a bound'' changes not only the fixed component but also the reduced gradient direction, which must be re-projected onto the null space of both $A$ and the newly active bounds. This is handled correctly in Algorithm~4 and the associated analysis.
  \item The projection onto the tangent space $T(x)$ requires solving a normal equations system whose structure changes as bounds become active or inactive. The block Cholesky factorization (Appendix~A), which reuses the constant factor $L$ of $AA^T$ and incrementally updates only the blocks involving the active bounds, is an efficient and well-engineered solution.
  \item The three-level iteration hierarchy (outer AL, inner trust region, minor projected CG) is carefully designed, with clear interfaces between levels and appropriate termination criteria at each.
\end{itemize}
That said, each individual ingredient---the AL framework, gradient projection with Cauchy point, projected CG, the normal equations approach---is well established in the literature. The novelty is one of \textbf{integration and engineering}: combining these pieces into a coherent algorithm for the polyhedral case, rather than introducing a fundamentally new algorithmic idea. The paper would benefit from being more explicit about this characterization.

\paragraph{Missing comparison.} A natural alternative to the polyhedral inner solver is \emph{variable elimination}: one could use the linear equalities $Ax = b$ to eliminate $m$ variables, reducing the problem to a bound-constrained one in $n - m$ variables, and then apply LANCELOT (or TRON) directly. The paper does not discuss why the direct polyhedral approach is preferable. Possible arguments include: avoiding the fill-in and conditioning issues that arise from the elimination, preserving sparsity structure in the Jacobians, and handling problems where $A$ is large and sparse more naturally. These should be stated.

\subsection{Contribution 2: Structured SR1 Update for the AL Hessian}

\paragraph{Claim.} The Hessian of the AL function is approximated by a structured SR1 update that exploits the sum-of-squares structure of both the objective and the penalty term, following the secant equation philosophy of Biggs (1977) and Dennis et al.\ (NL2SOL, 1981), adapted to the AL setting in the spirit of Li et al.\ (2002).

\paragraph{Assessment.} The structured secant equation
\[
B_{k+1}s_k = (J_{k+1} - J_k)^Tr_{k+1} + (C_{k+1} - C_k)^T\bar{\lambda}_{k+1}
\]
is obtained by imposing individual secant equations on each component Hessian $\nabla^2 r_i$ and $\nabla^2 c_i$, then summing with the appropriate weights. This is the same strategy as in Li et al.\ (2002), who applied it to the Lagrangian Hessian $\nabla^2_{xx}\mathcal{L}$ within an SQP framework. The adaptation to the AL Hessian is straightforward because the AL Hessian has the same sum-of-squares structure, with $\bar{\lambda}_i$ replacing $\lambda_i$ and the additional $\mu C^TC$ term absorbed into the first-order part $B_k^{GN}$.

The choice of SR1 over BFGS is well motivated: SR1 does not require a curvature condition $\inner{s_k}{y_k} > 0$, which may fail when the Hessian approximation is indefinite; and the trust-region CG framework naturally handles indefinite Hessians. The discussion comparing SR1 to BFGS/DFP in the constrained context (Section~2.3.3, paragraph on projected Hessian updates) is informative.

The observation about hybrid updates is noteworthy: the authors tested a switching strategy (GN approximation when the relative reduction is small, full SR1 otherwise) and found it \emph{counterproductive} for the AL Hessian, unlike the unconstrained or SQP cases where it is standard practice. The conjectured explanation---that the second-order terms involving the Lagrange multipliers cannot be neglected even for small-residual problems---is a useful empirical insight, though it remains without theoretical backing.

\paragraph{Novelty assessment.} This contribution is \textbf{incremental but meaningful}. The structured secant equation itself is not new (it is a direct transposition of the Biggs/NL2SOL idea to a structurally identical Hessian), and the SR1 formula is standard. What is new is: (a)~the specific application to the AL Hessian rather than the Lagrangian Hessian, (b)~the choice of SR1 in this context with a clear justification, and (c)~the empirical finding about hybrid updates. The paper should more clearly delineate what is genuinely new versus what is a direct transposition from the existing literature.

\subsection{Contribution 3: Global Convergence Analysis}

\paragraph{Claim.} The paper establishes global convergence of the method by adapting the trust-region analysis of Conn et al.\ (1988a) to the polyhedral constraint structure and by relying on the AL convergence framework of Conn et al.\ (1996b).

\paragraph{Assessment.} The convergence analysis has two parts of unequal completeness.

\subparagraph{Inner convergence (Section~3.1).} This is the stronger part of the theoretical contribution. The adaptation of Lemmas~3.1--3.4 to the polyhedral case requires genuine technical work: the projected gradient path involves reduced gradients that change at breakpoints (as bounds become active), and the bounds on the model decrease must account for the interaction between the linear equalities and the active bounds. The key steps---the non-increasing norm of the reduced gradient (Lemma~3.1), the lower bound on $\|z_k(t_k^{(1)})\|$ (Lemma~3.2), the model decrease bounds (Lemma~3.3), and the guaranteed decrease (Lemma~3.4)---are all correctly adapted, with the caveats noted in Sections~5.3--5.4 of this report. Once these are established, the convergence $h_k \to 0$ follows from the standard trust-region theory without further modification.

\subparagraph{Outer convergence (Section~3.2).} The global convergence theorem (Theorem~3.2) is \textbf{stated but not proved}. The paper asserts that it follows the framework of Conn et al.\ (1996b), but provides neither a proof nor a sketch. As discussed in Section~5.6 of this report, verifying that the hypotheses of the referenced results are satisfied in the present setting is non-trivial and constitutes an essential part of the claimed contribution. Without it, the paper cannot credibly claim to have ``established'' global convergence of the full algorithm.

\paragraph{Novelty assessment.} The inner convergence analysis is a \textbf{solid adaptation} of known theory to a new setting. The outer convergence result, if proved, would complete the picture but would not constitute a major theoretical advance---it would primarily confirm that the standard AL convergence framework applies. The convergence result is only first-order (no rate of convergence), which is standard for trust-region methods but limits the theoretical impact. A discussion of superlinear convergence under additional assumptions (e.g., Dennis--Mor\'e condition on the Hessian approximations) would strengthen the theoretical contribution.

\subsection{Industrial Motivation and Relationship to the Companion Preprint}

The TRAULLS algorithm is not developed in a vacuum. A companion preprint by the same first and second authors (Borie, Marcotte, Bastin, Dellacherie, ``Modernization of an optimization algorithm operated for short-term load forecasting in Quebec'') documents the first step of a broader modernization effort at Hydro-Qu\'ebec (HQ): the faithful conversion of the legacy ENLSIP solver from Fortran77 to Julia, published as the Enlsip.jl package (JOSS, 2024). That preprint, together with the JOSS paper, provides essential context for why the present algorithm is needed.

\paragraph{The scalability limitation.} The JOSS paper explicitly states that ENLSIP has ``no formal proof of convergence'' and that ``the algorithm is not suitable for large-scale applications,'' with performance degrading for $n + \ell \ge 200$ and $m \ge 1000$. Yet the operational calibration instances at HQ already reach 726 parameters and 14,040 residuals---well beyond ENLSIP's comfort zone. The preprint further notes that ``the lack of scalability of the underlying optimization method prevents from augmenting the model with additional features provided by upgraded weather stations, or to calibrate the parameters with larger datasets.'' This is the concrete industrial problem that TRAULLS is designed to solve: a trust-region augmented Lagrangian method that can handle polyhedral constraints, exploit least-squares structure, and scale to larger problem dimensions.

\paragraph{The incremental modernization strategy.} The preprint describes a deliberate two-phase approach: first convert the existing solver faithfully (Enlsip.jl), then develop an enhanced replacement (TRAULLS). This strategy reflects the operational constraint that ``the continuity of service is a hard constraint, and any modification must be introduced gradually.'' The TRAULLS draft is the second phase. This context strengthens the motivation for the present paper considerably---the algorithm is not merely an academic exercise but is intended for deployment in a critical industrial system.

\paragraph{Implications for the TRAULLS draft.} The existence of the operational test infrastructure (45 calibration instances with MAPE/RMSE metrics, Fortran--Julia interface, 10-minute time window constraint) means that the numerical experiments section of the TRAULLS draft could---and should---include comparisons against ENLSIP on these same instances, in addition to standard CUTEst benchmarks. This would provide the most compelling demonstration of the algorithm's practical value. The TRAULLS introduction should reference the companion preprint to make the industrial motivation explicit.

\subsection{Missing Elements}

Several elements, if included, would substantially strengthen the paper and its contribution claims.

\paragraph{Numerical experiments.} The numerical experiments section is empty (``TODO''). For an algorithmic paper, this is \textbf{critical}. Without computational evidence, the practical value of the approach remains undemonstrated. The experiments should include:
\begin{itemize}
  \item Comparisons against existing solvers: at minimum LANCELOT (bound-constrained AL), and ideally an SQP solver exploiting least-squares structure (e.g., the method of Li et al.\ (2002)) and a regularized approach (e.g., Bergou et al.\ (2021) or Orban and Siqueira (2020)).
  \item A standard test set: CUTEst problems with nonlinear least-squares structure and mixed constraints would be natural candidates.
  \item Ablation studies: the effect of the structured SR1 update versus the plain GN approximation, and the impact of the polyhedral solver versus variable elimination, should be quantified.
  \item Demonstration of the hybrid update observation: showing on concrete examples that the hybrid switching rule degrades performance for the AL Hessian would strengthen the empirical claim in Section~2.3.3.
\end{itemize}

\paragraph{Literature review.} As noted by the co-author's comments in the draft, the literature review needs expansion. In particular:
\begin{itemize}
  \item The positioning relative to Arreckx et al.\ (2016), who also implement an AL framework with slack variables and the magical step, should be sharpened. What does the present algorithm do that theirs does not (answer: exploit least-squares structure via the structured SR1 update; handle polyhedral constraints directly)?
  \item Recent developments in AL methods and nonlinear least-squares solvers should be surveyed more thoroughly.
\end{itemize}

\paragraph{Local convergence rate.} The paper establishes only global convergence to first-order critical points. For a structured quasi-Newton method, one would expect a discussion of whether superlinear convergence can be achieved. In the unconstrained NLS setting, Dennis et al.\ (1981) establish superlinear convergence of NL2SOL under appropriate conditions. Whether analogous results hold in the present AL framework is a natural question that, if addressed, would significantly increase the theoretical depth.

\subsection{Overall Characterization}

The paper's contributions are best characterized as a \textbf{careful synthesis} of established techniques---augmented Lagrangian methods, gradient projection with projected CG, structured quasi-Newton updates---into a coherent algorithm for a problem class (constrained NLS with polyhedral feasible sets) that has not been addressed by this particular combination before. The individual components are not new, but their integration is non-trivial and the engineering choices (SR1 over BFGS, normal equations with block Cholesky, the observation about hybrid updates) are sensible and well-argued.

The paper's value is primarily \textbf{practical and algorithmic} rather than theoretical. However, this case is currently undermined by two major gaps: the absence of numerical experiments and the missing outer convergence proof. Filling these would transform the paper from a solid working document into a complete and publishable contribution.


%% ============================================================
\section{Summary of Findings}\label{sec:summary}
%% ============================================================

\begin{table}[ht]
\centering
\renewcommand{\arraystretch}{1.25}
\small
\begin{tabular}{@{}p{1.5cm}p{6.2cm}p{1.8cm}p{3.5cm}@{}}
\toprule
\textbf{Location} & \textbf{Issue} & \textbf{Severity} & \textbf{Action Required} \\
\midrule
\multicolumn{4}{@{}l}{\textit{Mathematical correctness}} \\[2pt]
Eq.~(2.6) & Biconditional overstates stationarity & \statuswarn & Weaken to necessary cond. \\[3pt]
\S\,2.1 & Full rank of $\tilde{A}$ needs justification & \statuswarn & Add brief argument \\[3pt]
Appendix & $I_p - G_{21}G_{21}^T \succ 0$ not verified & \statuswarn & State Schur complement arg. \\[3pt]
Eq.~(2.33) & Typo: $c_{k+1}$ should be $x_{k+1}$ & \statuswarn & Fix typo \\[3pt]
Line~625 & $\Phi_K$ should be $\nabla_x\Phi_K$ & \statuswarn & Fix typo \\[3pt]
Lemma~3.2 & Chain of inequalities in wrong direction & \statuserr & Rewrite proof \\[3pt]
Lemma~3.4 & Unjustified step (TODO in draft) & \statuserr & Fill gap \\[3pt]
Thm.~3.2 & Proof entirely missing & \statuserr & Provide at least a sketch \\[3pt]
Line~928 & Subscript/superscript reversal & \statuswarn & Fix notation \\[3pt]
Eq.~(2.26) & Typo: $x_{x,j}$ should be $x_{k,j}$ & \statuswarn & Fix typo \\[6pt]
\multicolumn{4}{@{}l}{\textit{Contributions and completeness}} \\[2pt]
\S\,4 & Numerical experiments entirely missing & \statuserr & Add experiments and comparisons \\[3pt]
\S\,3.2 & Outer convergence proof missing & \statuserr & Provide proof or detailed sketch \\[3pt]
\S\,1 & Literature review incomplete & \statuswarn & Expand and sharpen positioning \\[3pt]
\S\,1 & No comparison with variable elimination & \statuswarn & Discuss alternative approaches \\[3pt]
\S\,2.3.3 & SR1 novelty vs.\ Li et al.\ unclear & \statuswarn & Delineate new vs.\ transposed \\[3pt]
\S\,3 & No local convergence rate discussed & \statuswarn & Add remark or defer explicitly \\[3pt]
\bottomrule
\end{tabular}
\caption{Summary of all issues identified in the mathematical developments and contribution claims.}
\label{tab:summary}
\end{table}

%% ============================================================
\section{Conclusion}
%% ============================================================

\paragraph{Mathematical validity.} The mathematical framework of the paper is fundamentally sound. The augmented Lagrangian formulation, the gradient and Hessian expressions, the projection machinery, the structured SR1 update, and the inner convergence theory are all correctly developed, with the caveats noted in Sections~\ref{sec:overall}--\ref{sec:typos}. The three items requiring the most attention are:

\begin{enumerate}[label=(\roman*)]
  \item \textbf{Lemma~3.2:} The proof is correct in substance but the presentation of the chain of inequalities should be rewritten for clarity.
  \item \textbf{Lemma~3.4:} The gap in Case~1 of the proof (the TODO note) should be filled by explicitly arguing that $\psi$ is strictly decreasing on $[0, t_k^{(5)}]$, implying $t_k^C \ge t_k^{(5)}$.
  \item \textbf{Theorem~3.2:} The missing proof of global convergence of the outer AL algorithm is the most significant gap. Even if the result follows the framework of Conn et al.\ (1996b), a sketch verifying that the hypotheses of the referenced results are satisfied in the present setting is essential.
\end{enumerate}

All other mathematical issues are minor (typos, notational inconsistencies, missing but straightforward justifications) and do not affect the validity of the main results.

\paragraph{Contributions.} The paper presents a careful synthesis of established techniques into a coherent algorithm for constrained NLS with polyhedral feasible sets. The integration is non-trivial and the design choices are well-motivated. The contributions are primarily practical and algorithmic in nature, and their case is currently undermined by two major gaps:

\begin{enumerate}[label=(\roman*)]
  \item \textbf{Numerical experiments:} The section is empty. Comparisons against existing solvers (LANCELOT, SQP methods, regularized approaches) on standard test sets are essential to demonstrate the practical value of the approach.
  \item \textbf{Outer convergence proof:} The missing proof prevents the paper from credibly claiming to have established global convergence of the full algorithm.
\end{enumerate}

\noindent Filling these two gaps, together with expanding the literature review and sharpening the positioning relative to the closest existing works, would transform the draft into a complete and publishable contribution.

\end{document}
